\clearpage
\setcounter{aufgabe}{22}
\einheitskopf{Einheit 4 von 5 · Grundwert berechnen}\label{e4}

\begin{aufgabe}{Grundwert – Was ist gegeben, was gesucht? Kreuze an, was in jedem Text gesucht ist. Rechne noch nicht.}
\begin{teile}
\teil 12 von 40 Schülern haben ein Haustier. Wie viel Prozent sind das? \\
\kreuz{Prozentwert}\kreuz{Prozentsatz}\kreuz{Grundwert}
\teil $35\,\%$ der 200 Karten sind verkauft. Wie viele Karten sind das? \\
\kreuz{Prozentwert}\kreuz{Prozentsatz}\kreuz{Grundwert}
\teil $18\,€$ sind $30\,\%$ des Preises. Wie hoch ist der Preis? \\
\kreuz{Prozentwert}\kreuz{Prozentsatz}\kreuz{Grundwert}
\teil In einer Klasse mit 28 Kindern fehlen $25\,\%$. Wie viele Kinder fehlen? \\
\kreuz{Prozentwert}\kreuz{Prozentsatz}\kreuz{Grundwert}
\teil 9 Plätze sind $15\,\%$ aller Plätze im Saal. Wie viele Plätze gibt es? \\
\kreuz{Prozentwert}\kreuz{Prozentsatz}\kreuz{Grundwert}
\end{teile}
\end{aufgabe}

\begin{aufgabe}{Grundwert – am Streifen. Der graue Teil ist gegeben, gesucht ist der ganze Streifen. Beispiel: Sind 2 der 10 Kästchen grau und $6\,€$ wert, so ist ein Kästchen $3\,€$ und der ganze Streifen $30\,€$.}
\streifenwertreihe{a/20/8 Stück/?, b/50/35 €/?, c/10/4 kg/?, d/25/12 m/?}
\begin{teilezwei}
\tz \leerfeld & \tz \leerfeld \\
\tz \leerfeld & \tz \leerfeld \\
\end{teilezwei}
\end{aufgabe}

\begin{aufgabe}{Grundwert – rechnen. Rechne mit dem Ein-Prozent-Weg: erst den Prozentwert durch den Prozentsatz, dann mal 100.}
\beispiel{\text{21\,€ sind 35\,\%:}\quad 1\,\% &= 21\,€ : 35 = 0{,}60\,€ \\ 100\,\% &= 100 \cdot 0{,}60\,€ = 60\,€}
\begin{teilezwei}
\tz $50\,\%$ sind $14\,€$: \leerfeld & \tz $25\,\%$ sind $9\,$kg: \leerfeld \\
\tz $20\,\%$ sind $7\,$m: \leerfeld & \tz $10\,\%$ sind $6\,€$: \leerfeld \\
\tz $40\,\%$ sind $18\,$kg: \leerfeld & \tz $12\,\%$ sind $42\,€$: \leerfeld \\
\end{teilezwei}
\end{aufgabe}

\begin{aufgabe}{Grundwert – aus dem Sachtext.}
\begin{teile}
\teil In einer Klasse sind 18 Kinder Mädchen. Das sind $60\,\%$ der Klasse. Wie viele Kinder hat die Klasse? \leerfeld
\teil Beim Kauf eines Druckers spart Herr Kunz $27\,€$. Das sind $15\,\%$ des alten Preises. Wie hoch war der alte Preis? \hfill (P10 2025) \\ \leerfeld[€]
\end{teile}
\end{aufgabe}

\begin{aufgabe}{Grundwert – Fehler finden. 35 Gäste eines Festes sind bereits da. Das sind $25\,\%$ aller Gäste. Tim will wissen, wie viele Gäste erwartet werden, und rechnet:}
\rechnung{0{,}25 \cdot 35 &= 8{,}75}
\begin{teile}
\teil Benenne den Fehler und korrigiere die Rechnung. \\ Fehler: \dsleer\dsleer\dsleer \quad richtig: \leerfeld
\teil 60 Karten sind $20\,\%$ aller Karten. Wie viele Karten gibt es? \leerfeld
\end{teile}
\end{aufgabe}

\begin{aufgabe}{Grundwert – Begründen. Begründe in einem Satz.}
\begin{teile}
\teil Warum nimmt man beim Ein-Prozent-Weg am Ende mit 100 mal? \\ \dsleer\dsleer\dsleer\dsleer
\teil Mia sagt: „Wenn der Prozentsatz kleiner als $100\,\%$ ist, dann ist der Grundwert größer als der Prozentwert.“ Hat sie recht? \janein \\ Begründung: \dsleer\dsleer\dsleer\dsleer
\end{teile}
\end{aufgabe}

\begin{aufgabe}{Gemischt. Entscheide bei jeder Aufgabe zuerst, ob Prozentwert, Prozentsatz oder Grundwert gesucht ist. Rechne erst danach.}
\begin{teile}
\teil In einem Korb liegen 40 Eier, $15\,\%$ davon sind zerbrochen. Wie viele Eier sind zerbrochen? \leerfeld
\teil Von 56 Plätzen sind 14 frei. Wie viel Prozent der Plätze sind frei? \leerfeld[\%]
\teil $33\,€$ sind $55\,\%$ des Preises. Wie hoch ist der Preis? \leerfeld[€]
\teil Ein Betrieb hat 120 Beschäftigte, $45\,\%$ arbeiten in Teilzeit. Wie viele sind das? \leerfeld
\teil 8 Kinder sind $16\,\%$ einer Gruppe. Wie viele Kinder hat die Gruppe? \leerfeld
\end{teile}
\end{aufgabe}
